% Source: https://artofproblemsolving.com/community/c3h2385070_counterintuitive_problem
% Stable identity: aops-topic-2385070-q01
Identical cubes with edge $1$ cm are stacked in a single column on a horizontal table: there is exactly one cube on each level, the bottom cube stands on the table, every other cube stands on the cube below it, and the cubes may be shifted sideways towards the edge of the table. The stack balances when, for every $k$, the centre of mass of the top $k$ cubes lies above the cube directly below them (above the table, for the whole stack); lying exactly above an edge is allowed. The \emph{overhang} is the distance by which the stack reaches beyond the edge of the table. Let $H_n=1+\dfrac12+\dfrac13+\cdots+\dfrac1n$.
\begin{enumerate}
\item[(a)] Prove that a balanced stack of $n$ cubes has overhang at most $\dfrac{H_n}{2}$ cm, and that this overhang can be reached.
\item[(b)] Deduce that an overhang of $1$ metre needs more than $2^{200}$ (which is more than $10^{60}$) cubes.
\end{enumerate}
